%%%%%%%%%%%%%%%%%%%%%%%%%%%%%%%%%%%%%%%%%%%%%%%%%%%%%%%%%%%%%%%%%%%
% Physik 4 - Blatt 11 Abgabe
%%%%%%%%%%%%%%%%%%%%%%%%%%%%%%%%%%%%%%%%%%%%%%%%%%%%%%%%%%%%%%%%%%%

\documentclass[11pt,a4paper]{scrartcl}

% Layout
\usepackage[margin=2.5cm]{geometry}
\usepackage[onehalfspacing]{setspace}

% General packages
\usepackage{graphicx}
\usepackage{enumitem}
\usepackage{tcolorbox}
\tcbuselibrary{breakable}
\usepackage[breaklinks=true,colorlinks=true,linkcolor=blue,urlcolor=blue,citecolor=blue]{hyperref}
\usepackage{amsmath,amsthm,amssymb,mathtools}
\usepackage{icomma}
\usepackage[T1]{fontenc}
\usepackage[utf8]{inputenc}
\usepackage[ngerman]{babel}
\usepackage{tikz}

% Units / tables / floats / references
\usepackage[locale = DE,space-before-unit=true,per-mode = symbol]{siunitx}
\usepackage{booktabs,multirow}
\usepackage{placeins}
\usepackage{csquotes}
\usepackage{pgfplots}
\pgfplotsset{compat=1.18}

% Optional math shortcuts
\newcommand{\R}{\mathbb{R}}
\newcommand{\N}{\mathbb{N}}
\newcommand{\Q}{\mathbb{Q}}
\newcommand{\set}[1]{\left\{ #1 \right\}}
\newcommand{\abs}[1]{\left| #1 \right|}
\newcommand{\norm}[1]{\left\lVert #1 \right\rVert}
\newcommand{\ol}[1]{\overline{#1}}
\newcommand{\half}{\frac{1}{2}}
\newcommand{\dd}{\,\mathrm{d}}
\newcommand{\hbarc}{\hbar c}

% Optional units
\DeclareSIUnit{\barn}{b}
\DeclareSIUnit{\nanobarn}{nb}

% Box styles
\newtcolorbox{calculationbox}[1][]{
    title=Calculation,
    breakable,
    #1
}

\newtcolorbox{solutionbox}[1][]{
    title=Solution,
    breakable,
    #1
}

\newtcolorbox{answerbox}[1][]{
    title=Answer,
    breakable,
    #1
}

\newtcolorbox{notebox}[1][]{
    title=Notes,
    breakable,
    #1
}

\newtcolorbox{sidecalcbox}[1][]{
    title=Side Calculation,
    breakable,
    #1
}

%%%%%%%%%%%%%%%%%%%%%%%%%%%%%%%%%%%%%%%%%%%%%%%%%%%%%%%%%%%%%%%%%%%
\begin{document}

    \title{Physik 4 - Blatt 11}
    \author{Tobias Laser}
    \date{\today}
    \maketitle
    \vfill
    \thispagestyle{empty}
    \newpage

    \pagenumbering{arabic}

    \paragraph{Aufgabe 1: $Z$-Resonanz}

    Am Large-Electron-Positron Collider (LEP) des CERN wurden von 1989 bis 1995
    Elektronen und Positronen so weit beschleunigt und zur Kollision gebracht, dass
    im Schwerpunktsystem die Ruheenergie des $Z$-Bosons zur Verfügung stand:
    \[
        E_{\mathrm{CM}}=m_Zc^2=\SI{91.2}{GeV}.
    \]
    Bei dieser Energie ist der Wirkungsquerschnitt für die Produktion des $Z$-Bosons maximal.
    Der differentielle Wirkungsquerschnitt sei
    \[
        \frac{\dd\sigma}{\dd\Omega}(E_{\mathrm{CM}},\theta)
        = (\hbar c)^2\frac{\alpha^2}{s}\frac{G(s)}{4}
        \left(1+\cos^2\theta\right),
    \]
    wobei
    \[
        G(s)=1+\frac{s^2g_eg_\mu}{(s-m_Z^2c^4)^2+m_Z^2c^4\Gamma_Z^2}.
    \]
    Außerdem gilt
    \[
        s=E_{\mathrm{CM}}^2,\qquad
        \Gamma_Z=\SI{2.5}{GeV},\qquad
        g_e=g_\mu=0.355.
    \]
    Der erste Term in $G(s)$ beschreibt den Beitrag über ein virtuelles Photon,
    der zweite Term den Beitrag über das $Z$-Boson.

    \begin{enumerate}[label=(\alph*)]
        \item
            Um welchen Faktor ist der zum $Z$-Boson gehörige Term am Maximum der Resonanz
            größer als der Beitrag des virtuellen Photons?

            \begin{calculationbox}
                Am Maximum der Resonanz gilt
                \[
                    E_{\mathrm{CM}}=m_Zc^2
                    \qquad\Longrightarrow\qquad
                    s=m_Z^2c^4.
                \]
                Damit verschwindet im Nenner des $Z$-Terms der Ausdruck
                $(s-m_Z^2c^4)^2$.

                \begin{calculationbox}
                    Der Photon-Beitrag in $G(s)$ ist einfach
                    \[
                        G_\gamma=1.
                    \]
                    Der $Z$-Beitrag am Resonanzmaximum ist
                    \begin{align*}
                        G_Z
                        &=\frac{s^2g_eg_\mu}{m_Z^2c^4\Gamma_Z^2}.
                    \end{align*}
                    Da am Maximum $s=m_Z^2c^4$ gilt, wird daraus
                    \[
                        G_Z=\frac{s g_eg_\mu}{\Gamma_Z^2}.
                    \]
                \end{calculationbox}

                \begin{calculationbox}
                    Mit
                    \[
                        s=E_{\mathrm{CM}}^2=(\SI{91.2}{GeV})^2
                    \]
                    folgt
                    \begin{align*}
                        G_Z
                        &=\frac{(91.2)^2\cdot 0.355\cdot 0.355}{(2.5)^2}\\
                        &=167.7.
                    \end{align*}
                    Der gesuchte Faktor gegenüber dem virtuellen Photon ist daher
                    \[
                        \frac{G_Z}{G_\gamma}=167.7.
                    \]
                \end{calculationbox}
            \end{calculationbox}

            \begin{answerbox}
                Am Resonanzmaximum ist der $Z$-Beitrag etwa
                \[
                    \boxed{1.68\cdot 10^2}
                \]
                mal größer als der Beitrag des virtuellen Photons.
            \end{answerbox}

        \item
            Skizzieren Sie die Winkelabhängigkeit des differentiellen Wirkungsquerschnitts.
            Vergleichen Sie mit der entsprechenden Abhängigkeit des Rutherfordstreuquerschnitts.

            \begin{calculationbox}
                Für festes $E_{\mathrm{CM}}$ sind alle Faktoren bis auf den Winkel konstant.
                Daher ist die Winkelabhängigkeit
                \[
                    \frac{\dd\sigma}{\dd\Omega}\propto 1+\cos^2\theta.
                \]

                \begin{calculationbox}
                    Die Funktion ist symmetrisch unter
                    \[
                        \theta\mapsto \pi-\theta.
                    \]
                    Sie ist maximal bei
                    \[
                        \theta=0,\pi,
                    \]
                    denn dort ist $\cos^2\theta=1$ und damit
                    \[
                        1+\cos^2\theta=2.
                    \]
                    Sie ist minimal bei
                    \[
                        \theta=\frac{\pi}{2},
                    \]
                    denn dort ist $\cos^2\theta=0$ und damit
                    \[
                        1+\cos^2\theta=1.
                    \]
                \end{calculationbox}

                \begin{calculationbox}
                    Eine qualitative Skizze ist:
                    \begin{center}
                    \begin{tikzpicture}
                        \begin{axis}[
                            width=0.82\textwidth,
                            height=6cm,
                            xlabel={$\theta$},
                            ylabel={relative Stärke},
                            xmin=0, xmax=180,
                            ymin=0, ymax=2.25,
                            xtick={0,45,90,135,180},
                            ytick={0,1,2},
                            grid=both,
                            legend pos=south east
                        ]
                            \addplot[thick,domain=0:180,samples=200]{1+cos(x)^2};
                            \addlegendentry{$1+\cos^2\theta$}
                        \end{axis}
                    \end{tikzpicture}
                    \end{center}
                \end{calculationbox}

                \begin{calculationbox}
                    Beim Rutherfordstreuquerschnitt gilt dagegen qualitativ
                    \[
                        \frac{\dd\sigma_R}{\dd\Omega}\propto \frac{1}{\sin^4(\theta/2)}.
                    \]
                    Diese Winkelabhängigkeit ist stark nach vorne gerichtet und divergiert
                    für $\theta\to0$. Im Gegensatz dazu ist $1+\cos^2\theta$ überall endlich
                    und hat nur einen Faktor $2$ zwischen Minimum und Maximum.
                \end{calculationbox}
            \end{calculationbox}

            \begin{answerbox}
                Die Winkelabhängigkeit der $e^+e^-\to \mu^+\mu^-$-Reaktion ist glatt,
                endlich und vorwärts-rückwärts-symmetrisch:
                \[
                    \boxed{\frac{\dd\sigma}{\dd\Omega}\propto 1+\cos^2\theta.}
                \]
                Rutherfordstreuung ist dagegen stark vorwärtsdominant und besitzt eine
                Divergenz für kleine Streuwinkel.
            \end{answerbox}

        \item
            Berechnen Sie den totalen Wirkungsquerschnitt $\sigma_{\mathrm{total}}$,
            indem Sie den differentiellen Wirkungsquerschnitt über den gesamten Raumwinkel integrieren.

            \begin{calculationbox}
                Der totale Wirkungsquerschnitt ist
                \[
                    \sigma_{\mathrm{total}}
                    =\int \frac{\dd\sigma}{\dd\Omega}\,\dd\Omega.
                \]
                Wegen
                \[
                    \dd\Omega=\sin\theta\,\dd\theta\,\dd\varphi
                \]
                muss über $\varphi\in[0,2\pi]$ und $\theta\in[0,\pi]$ integriert werden.

                \begin{calculationbox}
                    Einsetzen liefert
                    \begin{align*}
                        \sigma_{\mathrm{total}}
                        &=(\hbar c)^2\frac{\alpha^2}{s}\frac{G(s)}{4}
                        \int_0^{2\pi}\int_0^\pi
                        (1+\cos^2\theta)\sin\theta\,\dd\theta\,\dd\varphi.
                    \end{align*}
                \end{calculationbox}

                \begin{calculationbox}
                    Zuerst integrieren wir über $\varphi$:
                    \[
                        \int_0^{2\pi}\dd\varphi=2\pi.
                    \]
                    Für das $\theta$-Integral setzen wir
                    \[
                        u=\cos\theta,
                        \qquad
                        \dd u=-\sin\theta\,\dd\theta.
                    \]
                    Dann gilt
                    \begin{align*}
                        \int_0^\pi (1+\cos^2\theta)\sin\theta\,\dd\theta
                        &=\int_{1}^{-1}(1+u^2)(-\dd u)\\
                        &=\int_{-1}^{1}(1+u^2)\,\dd u\\
                        &=\left[u+\frac{u^3}{3}\right]_{-1}^{1}\\
                        &=\frac{8}{3}.
                    \end{align*}
                \end{calculationbox}

                \begin{calculationbox}
                    Insgesamt ist also
                    \[
                        \int (1+\cos^2\theta)\,\dd\Omega
                        =2\pi\cdot\frac{8}{3}=\frac{16\pi}{3}.
                    \]
                    Damit folgt
                    \begin{align*}
                        \sigma_{\mathrm{total}}
                        &=(\hbar c)^2\frac{\alpha^2}{s}\frac{G(s)}{4}\cdot\frac{16\pi}{3}\\
                        &=(\hbar c)^2\frac{\alpha^2}{s}G(s)\frac{4\pi}{3}.
                    \end{align*}
                \end{calculationbox}
            \end{calculationbox}

            \begin{answerbox}
                Der totale Wirkungsquerschnitt ist
                \[
                    \boxed{
                    \sigma_{\mathrm{total}}(s)
                    =(\hbar c)^2\frac{\alpha^2}{s}G(s)\frac{4\pi}{3}.}
                \]
            \end{answerbox}

        \item
            Wie groß ist somit der totale Wirkungsquerschnitt bei $E_{\mathrm{CM}}=m_Zc^2$
            in Einheiten von Nanobarn? Welchen Radius hätte ein Kreis mit gleicher Fläche?
            Nutzen Sie $\hbar c=\SI{197}{MeV.fm}$.

            \begin{calculationbox}
                Wir verwenden
                \[
                    \hbar c=\SI{197}{MeV.fm}=\SI{0.197}{GeV.fm},
                    \qquad
                    \alpha=\frac{1}{137.036}.
                \]
                Aus Teil (a) haben wir am Maximum
                \[
                    G(s)=1+167.7=168.7.
                \]

                \begin{calculationbox}
                    Bei $E_{\mathrm{CM}}=\SI{91.2}{GeV}$ ist
                    \[
                        s=(\SI{91.2}{GeV})^2=\SI{8317.44}{GeV^2}.
                    \]
                    Daher
                    \begin{align*}
                        \sigma_{\mathrm{total}}
                        &=(0.197)^2\,\si{fm^2}\frac{(1/137.036)^2}{8317.44}
                        \cdot168.7\cdot\frac{4\pi}{3}\\
                        &=1.756\cdot10^{-7}\,\si{fm^2}.
                    \end{align*}
                \end{calculationbox}

                \begin{calculationbox}
                    Für die Umrechnung benutzen wir
                    \[
                        1\,\si{fm^2}=10^{-30}\,\si{m^2},
                        \qquad
                        1\,\si{barn}=10^{-28}\,\si{m^2}=100\,\si{fm^2}.
                    \]
                    Außerdem gilt
                    \[
                        1\,\si{barn}=10^9\,\si{nanobarn}.
                    \]
                    Also ist
                    \[
                        1\,\si{fm^2}=10^7\,\si{nanobarn}.
                    \]
                    Damit
                    \[
                        \sigma_{\mathrm{total}}
                        =1.756\cdot10^{-7}\,\si{fm^2}\cdot10^7\,\frac{\si{nanobarn}}{\si{fm^2}}
                        =1.756\,\si{nanobarn}.
                    \]
                \end{calculationbox}

                \begin{calculationbox}
                    Ein Kreis mit gleicher Fläche erfüllt
                    \[
                        \sigma_{\mathrm{total}}=\pi r^2.
                    \]
                    Also
                    \begin{align*}
                        r&=\sqrt{\frac{\sigma_{\mathrm{total}}}{\pi}}\\
                        &=\sqrt{\frac{1.756\cdot10^{-7}\,\si{fm^2}}{\pi}}\\
                        &=2.36\cdot10^{-4}\,\si{fm}\\
                        &=2.36\cdot10^{-19}\,\si{m}.
                    \end{align*}
                \end{calculationbox}
            \end{calculationbox}

            \begin{answerbox}
                Am Resonanzmaximum ergibt sich
                \[
                    \boxed{\sigma_{\mathrm{total}}=1.76\,\si{nanobarn}.}
                \]
                Ein Kreis mit gleicher Fläche hätte den Radius
                \[
                    \boxed{r=2.36\cdot10^{-4}\,\si{fm}=2.36\cdot10^{-19}\,\si{m}.}
                \]
            \end{answerbox}

        \item
            Plotten Sie $\sigma_{\mathrm{total}}$ als Funktion von $E_{\mathrm{CM}}$ im Bereich
            $50-150\,\si{GeV}$ in linearer Darstellung und im Bereich $10-500\,\si{GeV}$
            in doppeltlogarithmischer Darstellung.

            \begin{calculationbox}
                Für die Plots verwenden wir die Formel
                \[
                    \sigma_{\mathrm{total}}(E)
                    =(\hbar c)^2\frac{\alpha^2}{E^2}G(E^2)\frac{4\pi}{3}.
                \]
                Dabei wird das Ergebnis in Nanobarn dargestellt. Die verwendete Funktion ist
                \[
                    \sigma_{\mathrm{nb}}(E)
                    =10^7\,(0.197)^2\frac{(1/137.036)^2}{E^2}\frac{4\pi}{3}
                    \left[1+\frac{E^4\cdot0.355^2}{(E^2-91.2^2)^2+91.2^2\cdot2.5^2}\right].
                \]

                \begin{calculationbox}
                    Lineare Darstellung um die Resonanz:
                    \begin{center}
                    \begin{tikzpicture}
                        \begin{axis}[
                            width=0.86\textwidth,
                            height=6.5cm,
                            xlabel={$E_{\mathrm{CM}}$ [GeV]},
                            ylabel={$\sigma_{\mathrm{total}}$ [nb]},
                            xmin=50, xmax=150,
                            ymin=0, ymax=2.0,
                            grid=both
                        ]
                            \addplot[thick,domain=50:150,samples=400]
                            {1e7*(0.197^2)*(1/137.036)^2/(x^2)*(4*pi/3)*(1 + (x^4*0.355^2)/(((x^2-91.2^2)^2)+(91.2^2*2.5^2)))};
                        \end{axis}
                    \end{tikzpicture}
                    \end{center}
                \end{calculationbox}

                \begin{calculationbox}
                    Doppeltlogarithmische Darstellung im größeren Energiebereich:
                    \begin{center}
                    \begin{tikzpicture}
                        \begin{loglogaxis}[
                            width=0.86\textwidth,
                            height=6.5cm,
                            xlabel={$E_{\mathrm{CM}}$ [GeV]},
                            ylabel={$\sigma_{\mathrm{total}}$ [nb]},
                            xmin=10, xmax=500,
                            ymin=1e-4, ymax=3,
                            grid=both
                        ]
                            \addplot[thick,domain=10:500,samples=500]
                            {1e7*(0.197^2)*(1/137.036)^2/(x^2)*(4*pi/3)*(1 + (x^4*0.355^2)/(((x^2-91.2^2)^2)+(91.2^2*2.5^2)))};
                        \end{loglogaxis}
                    \end{tikzpicture}
                    \end{center}
                \end{calculationbox}
            \end{calculationbox}

            \begin{answerbox}
                Der lineare Plot zeigt den scharfen Resonanzpeak bei
                \[
                    E_{\mathrm{CM}}\approx\SI{91.2}{GeV}.
                \]
                In der doppeltlogarithmischen Darstellung sieht man zusätzlich den breiteren
                Verlauf außerhalb der Resonanz.
            \end{answerbox}
    \end{enumerate}

    \newpage
    \paragraph{Aufgabe 2: Solare Neutrinos}

    Die Netto-Reaktion, mit der die Sonne Energie erzeugt, lautet
    \[
        4\,{}^1\mathrm{H}\rightarrow{}^4\mathrm{He}+2e^++2e^-+2\nu_e.
    \]
    Berücksichtigt man, dass die entstehenden Positronen mit Elektronen annihilieren,
    wird pro Netto-Reaktion eine Energie von
    \[
        Q=\SI{26.73}{MeV}
    \]
    frei. Pro Reaktion entstehen zwei Elektron-Neutrinos.

    \begin{enumerate}[label=(\alph*)]
        \item
            Die Solarkonstante beträgt
            \[
                S_0=\SI{1361}{W/m^2}.
            \]
            Berechnen Sie daraus den Fluss der Neutrinos im Abstand der Erde.

            \begin{calculationbox}
                Die Solarkonstante ist eine Leistung pro Fläche. Daher gibt
                \[
                    \frac{S_0}{Q}
                \]
                die Anzahl der Fusionsreaktionen pro Sekunde und Quadratmeter an,
                deren Energiefluss bei der Erde ankommt. Da pro Reaktion zwei Neutrinos
                entstehen, ist der Neutrinofluss
                \[
                    \Phi_\nu=2\frac{S_0}{Q}.
                \]

                \begin{calculationbox}
                    Zunächst rechnen wir die freiwerdende Energie in Joule um:
                    \begin{align*}
                        Q&=\SI{26.73}{MeV}\\
                        &=26.73\cdot10^6\,\si{eV}\cdot1.602\cdot10^{-19}\,\frac{\si{J}}{\si{eV}}\\
                        &=4.283\cdot10^{-12}\,\si{J}.
                    \end{align*}
                \end{calculationbox}

                \begin{calculationbox}
                    Damit folgt
                    \begin{align*}
                        \Phi_\nu
                        &=2\frac{1361\,\si{J.s^{-1}.m^{-2}}}{4.283\cdot10^{-12}\,\si{J}}\\
                        &=6.36\cdot10^{14}\,\si{m^{-2}.s^{-1}}.
                    \end{align*}
                \end{calculationbox}
            \end{calculationbox}

            \begin{answerbox}
                Der solare Neutrinofluss an der Erde beträgt näherungsweise
                \[
                    \boxed{\Phi_\nu=6.36\cdot10^{14}\,\si{m^{-2}.s^{-1}}.}
                \]
            \end{answerbox}

        \item
            Wie viele Neutrinos gehen pro Sekunde durch Ihren Daumennagel, wenn dessen Fläche
            \(\SI{1}{cm^2}\) beträgt?

            \begin{calculationbox}
                Die Fläche ist
                \[
                    A=\SI{1}{cm^2}=10^{-4}\,\si{m^2}.
                \]
                Die Anzahl pro Sekunde ist
                \[
                    \dot N=\Phi_\nu A.
                \]

                \begin{calculationbox}
                    Einsetzen liefert
                    \begin{align*}
                        \dot N
                        &=6.36\cdot10^{14}\,\si{m^{-2}.s^{-1}}\cdot10^{-4}\,\si{m^2}\\
                        &=6.36\cdot10^{10}\,\si{s^{-1}}.
                    \end{align*}
                \end{calculationbox}
            \end{calculationbox}

            \begin{answerbox}
                Pro Sekunde gehen durch einen Daumennagel mit Fläche \(\SI{1}{cm^2}\) etwa
                \[
                    \boxed{6.36\cdot10^{10}}
                \]
                solare Neutrinos.
            \end{answerbox}

        \item
            Die solaren Neutrinos haben eine typische Energie von \(\SI{100}{keV}\).
            In diesem Energiebereich sei der Wirkungsquerschnitt für Reaktionen von
            Neutrinos und Nukleonen
            \[
                \sigma(\nu_eN)\simeq10^{-21}\,\si{barn}.
            \]
            Berechnen Sie damit die Reaktionsrate bei einer Person mit einer Masse von
            \(\SI{70}{kg}\), sowohl pro Sekunde als auch pro Jahr.

            \begin{calculationbox}
                Die Reaktionsrate ist näherungsweise
                \[
                    R=\Phi_\nu\sigma N_N,
                \]
                wobei $N_N$ die Anzahl der Nukleonen im Körper ist.

                \begin{calculationbox}
                    Der Wirkungsquerschnitt in SI-Einheiten ist
                    \[
                        \sigma=10^{-21}\,\si{barn}
                        =10^{-21}\cdot10^{-28}\,\si{m^2}
                        =10^{-49}\,\si{m^2}.
                    \]
                \end{calculationbox}

                \begin{calculationbox}
                    Die Nukleonenzahl einer Person mit Masse $m=\SI{70}{kg}$ schätzen wir über
                    die Protonenmasse ab:
                    \[
                        m_N\approx1.673\cdot10^{-27}\,\si{kg}.
                    \]
                    Damit
                    \begin{align*}
                        N_N&\approx\frac{70}{1.673\cdot10^{-27}}\\
                        &=4.19\cdot10^{28}.
                    \end{align*}
                \end{calculationbox}

                \begin{calculationbox}
                    Damit ist die Reaktionsrate
                    \begin{align*}
                        R&=6.36\cdot10^{14}\,\si{m^{-2}.s^{-1}}
                        \cdot10^{-49}\,\si{m^2}\cdot4.19\cdot10^{28}\\
                        &=2.66\cdot10^{-6}\,\si{s^{-1}}.
                    \end{align*}
                    Pro Jahr mit
                    \[
                        1\,\si{a}=365.25\cdot24\cdot3600\,\si{s}=3.156\cdot10^7\,\si{s}
                    \]
                    folgt
                    \begin{align*}
                        N_{\mathrm{Jahr}}&=R\cdot3.156\cdot10^7\\
                        &=83.9.
                    \end{align*}
                \end{calculationbox}
            \end{calculationbox}

            \begin{answerbox}
                Für eine Person mit \(\SI{70}{kg}\) ergibt sich
                \[
                    \boxed{R=2.66\cdot10^{-6}\,\si{s^{-1}}.}
                \]
                Das entspricht etwa
                \[
                    \boxed{84\text{ Reaktionen pro Jahr}.}
                \]
            \end{answerbox}

        \item
            Eine solche Reaktion sehe vereinfacht so aus:
            \[
                \nu_e+N\rightarrow e^-+N'.
            \]
            Das freiwerdende Elektron verursacht Strahlenschäden im menschlichen Gewebe.
            Berechnen Sie, wie viel Energie durch Elektronen im Körper einer Person in einem Jahr
            frei wird, wenn im Mittel die Hälfte der Energie der Neutrinos auf die Elektronen übergeht.

            \begin{calculationbox}
                Die typische Neutrinoenergie ist
                \[
                    E_\nu=\SI{100}{keV}.
                \]
                Im Mittel wird davon die Hälfte auf das Elektron übertragen:
                \[
                    E_e=\frac{1}{2}E_\nu=\SI{50}{keV}.
                \]

                \begin{calculationbox}
                    In Joule ist das
                    \begin{align*}
                        E_e&=50\cdot10^3\,\si{eV}\cdot1.602\cdot10^{-19}\,\frac{\si{J}}{\si{eV}}\\
                        &=8.01\cdot10^{-15}\,\si{J}.
                    \end{align*}
                \end{calculationbox}

                \begin{calculationbox}
                    Aus Teil (c) hatten wir ungefähr
                    \[
                        N_{\mathrm{Jahr}}=83.9
                    \]
                    Reaktionen pro Jahr. Daher ist die pro Jahr im Körper abgegebene Energie
                    \begin{align*}
                        \Delta E_{\mathrm{Jahr}}
                        &=N_{\mathrm{Jahr}}E_e\\
                        &=83.9\cdot8.01\cdot10^{-15}\,\si{J}\\
                        &=6.72\cdot10^{-13}\,\si{J}.
                    \end{align*}
                \end{calculationbox}
            \end{calculationbox}

            \begin{answerbox}
                Durch diese Elektronen wird pro Jahr ungefähr die Energie
                \[
                    \boxed{\Delta E_{\mathrm{Jahr}}=6.72\cdot10^{-13}\,\si{J}}
                \]
                im Körper abgegeben.
            \end{answerbox}

        \item
            Die Äquivalentdosis sei
            \[
                H_\nu=\left(\frac{\Delta E}{m}\right)w_R,
            \]
            wobei für Elektronen $w_R=1$ gilt. Berechnen Sie die jährliche Äquivalentdosis
            durch solare Neutrinos und vergleichen Sie diese mit der natürlichen Strahlenbelastung
            im Bereich von \(\SI{1}{mSv/a}\). Dabei gilt \(\SI{1}{Sv}=\SI{1}{J/kg}\).

            \begin{calculationbox}
                Wir verwenden
                \[
                    \Delta E=6.72\cdot10^{-13}\,\si{J},
                    \qquad
                    m=\SI{70}{kg},
                    \qquad
                    w_R=1.
                \]
                Damit ist
                \[
                    H_\nu=\frac{\Delta E}{m}.
                \]

                \begin{calculationbox}
                    Einsetzen ergibt
                    \begin{align*}
                        H_\nu
                        &=\frac{6.72\cdot10^{-13}\,\si{J}}{70\,\si{kg}}\\
                        &=9.61\cdot10^{-15}\,\si{J/kg}\\
                        &=9.61\cdot10^{-15}\,\si{Sv/a}.
                    \end{align*}
                \end{calculationbox}

                \begin{calculationbox}
                    In Millisievert ist das
                    \[
                        H_\nu=9.61\cdot10^{-12}\,\si{mSv/a}.
                    \]
                    Verglichen mit
                    \[
                        H_{\mathrm{nat}}\sim\SI{1}{mSv/a}
                    \]
                    ist der Anteil
                    \begin{align*}
                        \frac{H_\nu}{H_{\mathrm{nat}}}
                        &=\frac{9.61\cdot10^{-12}\,\si{mSv/a}}{1\,\si{mSv/a}}\\
                        &=9.61\cdot10^{-12}.
                    \end{align*}
                    Die Dosis durch solare Neutrinos ist also etwa
                    \[
                        \frac{1}{9.61\cdot10^{-12}}\approx1.0\cdot10^{11}
                    \]
                    mal kleiner als die natürliche jährliche Strahlenbelastung.
                \end{calculationbox}
            \end{calculationbox}

            \begin{answerbox}
                Die jährliche Äquivalentdosis durch solare Neutrinos beträgt
                \[
                    \boxed{H_\nu=9.61\cdot10^{-15}\,\si{Sv/a}
                    =9.61\cdot10^{-12}\,\si{mSv/a}.}
                \]
                Das ist ungefähr
                \[
                    \boxed{10^{11}\text{-mal kleiner}}
                \]
                als eine natürliche Strahlenbelastung von etwa \(\SI{1}{mSv/a}\).
            \end{answerbox}
    \end{enumerate}

\end{document}
